\begin{figure}[H]
  \centering
  \resizebox{0.98\linewidth}{!}{%
  \begin{tikzpicture}[
    font=\small,>=Latex,
    box/.style={draw,rounded corners=1.2pt,align=center,minimum height=8mm,inner sep=4pt},
    data/.style={box,fill=blue!7}, ctrl/.style={box,fill=orange!10},
    use/.style={box,fill=green!9}, arr/.style={->,thick}, dasharr/.style={->,thick,dashed}]
    \node[data] (addr) {IDU 地址基址选择\\GPR / LSU 快速结果 / WBU 普通结果};
    \node[data,right=13mm of addr] (add) {22-bit 地址加法\\专用 ID/EX payload};
    \node[ctrl,right=13mm of add] (index) {注册 query index\\tag 在 EXU 比较};
    \node[data,right=13mm of index] (bram) {主数据 BRAM\\C1 同步读};
    \node[use,right=13mm of bram] (lsu) {LSU 对齐与扩展\\命中结果前递};
    \node[ctrl,below=14mm of add] (blocked) {不进入地址锥\\D-cache 数据 / 长算术 / 加速 lane};
    \node[ctrl,below=14mm of bram] (consumers) {依赖消费者\\数据未有效时按 RAW 停顿};
    \draw[arr] (addr)--(add);
    \draw[arr] (add)--(index);
    \draw[arr] (index)--(bram);
    \draw[arr] (bram)--(lsu);
    \draw[dasharr] (blocked)--node[below]{隔离}(add);
    \draw[arr] (lsu)|-(consumers);
  \end{tikzpicture}}
  \caption{同步 load 通路与地址生成隔离边界}\label{fig:late-load}
\end{figure}
